\documentclass[aspectratio=169,11pt]{beamer}
\usetheme{TML}

\aula{2}{Introdução a Python, NumPy e JAX}

\newcommand{\nbdir}{aulas/02-python-numpy-jax/notebooks}

\begin{document}

\maketitle

% =====================================================================
\begin{frame}{Por que um economista deveria programar?}
  \begin{columns}[T]
    \begin{column}{0.5\textwidth}
      \textbf{Uma tarefa comum:} calcular a inflação acumulada em 12 meses a
      partir do IPCA mensal.
      \begin{itemize}
        \item Na planilha: arrastar fórmulas, copiar, colar\ldots
        \item E quando chegar o dado do mês que vem?
        \item E para 27 UFs? Para 400 subitens?
      \end{itemize}
    \end{column}
    \begin{column}{0.46\textwidth}
      \textbf{Com código:}
      \begin{itemize}
        \item \alert{Reprodutível}: alguém roda e obtém o mesmo número
        \item \alert{Escalável}: 1 série ou 10 mil, mesmo esforço
        \item \alert{Auditável}: cada decisão está escrita
      \end{itemize}
      \medskip
      Na Aula 1, o \py{fit} e o \py{predict} rodaram em Python. Hoje
      aprendemos a linguagem por trás deles.
    \end{column}
  \end{columns}
\end{frame}

\begin{frame}{Roteiro da aula}
  \setbeamertemplate{section in toc}[sections numbered]
  \tableofcontents[hideallsubsections]
  \vfill
  \small\textcolor{gray}{1h de fundamentos \textbullet{} 30 min de demonstração \textbullet{} 30 min de laboratório}
\end{frame}

% =====================================================================
\section{O ambiente de trabalho}

\begin{frame}{Google Colab e notebooks Jupyter}
  \begin{columns}[c]
    \begin{column}{0.55\textwidth}
      \begin{itemize}
        \item Um \alert{notebook} intercala células de \textbf{código} e de
              \textbf{texto} (Markdown, com \LaTeX)
        \item O \alert{kernel} guarda o estado: variáveis vivem entre células
        \item O Colab roda no navegador, sem instalar nada, com GPU gratuita
        \item \texttt{Shift+Enter} executa a célula e avança
      \end{itemize}
    \end{column}
    \begin{column}{0.4\textwidth}
      \begin{atencao}
        \small O que importa é a \textbf{ordem em que você executou} as células,
        não a ordem em que elas aparecem. Na dúvida:
        \emph{Ambiente de execução $\to$ Reiniciar e executar tudo}.
      \end{atencao}
    \end{column}
  \end{columns}
\end{frame}

% =====================================================================
\section{Python essencial}

\begin{frame}[fragile]{Variáveis e tipos}
  Uma variável é um \alert{nome} associado a um \alert{valor}. O tipo é inferido.
  \begin{codigo}
    selic = 15.0                 # float: número real
    n_meses = 12                 # int: inteiro
    pais = "Brasil"              # str: texto
    acima_da_meta = True         # bool: verdadeiro/falso

    type(selic)                  # <class 'float'>
    taxa_mensal = (1 + selic / 100) ** (1 / n_meses) - 1
    print(f"Selic mensal equivalente: {taxa_mensal:.4%}")
  \end{codigo}
  \begin{atencao}
    \small Ponto flutuante não é exato: \py{0.1 + 0.2 == 0.3} é \py{False}.
    Compare números reais com tolerância (\py{abs(a - b) < 1e-9} ou \py{math.isclose}).
  \end{atencao}
\end{frame}

\begin{frame}[fragile]{Errar faz parte: como ler uma mensagem de erro}
  \begin{columns}[T]
    \begin{column}{0.5\textwidth}
      \begin{codigo}
        inflacao = "5"   # veio como texto
        inflacao + 3
        # TypeError: can only concatenate
        #   str (not "int") to str

        float(inflacao) + 3   # 8.0
      \end{codigo}
      \small Leia a \alert{última linha} primeiro: ela diz o \textbf{tipo} do
      erro e o motivo. As linhas acima mostram onde ele aconteceu.
    \end{column}
    \begin{column}{0.46\textwidth}
      \small
      \begin{tabular}{@{}l>{\raggedright\arraybackslash}p{0.6\linewidth}@{}}
        \toprule
        Erro & Causa mais comum \\
        \midrule
        \py{NameError}        & nome digitado errado ou célula não executada \\
        \py{TypeError}        & operação com o tipo errado (texto + número) \\
        \py{IndexError}       & índice além do fim da lista \\
        \py{IndentationError} & recuo inconsistente num bloco \\
        \bottomrule
      \end{tabular}
    \end{column}
  \end{columns}
\end{frame}

\begin{frame}[fragile]{Coleções: listas e dicionários}
  \begin{columns}[T]
    \begin{column}{0.48\textwidth}
      \textbf{Lista}: sequência ordenada
      \begin{codigo}
        ipca = [0.16, 0.43, 0.56, 0.43]
        ipca[0]     # 0.16 (índice 0!)
        ipca[-1]    # 0.43 (último)
        ipca[1:3]   # [0.43, 0.56]
        ipca.append(0.26)
        len(ipca)   # 5
      \end{codigo}
    \end{column}
    \begin{column}{0.48\textwidth}
      \textbf{Dicionário}: pares chave $\to$ valor
      \begin{codigo}
        metas = {
            "centro": 3.0,
            "tolerancia": 1.5,
        }
        metas["centro"]   # 3.0
        teto = (metas["centro"]
                + metas["tolerancia"])
      \end{codigo}
    \end{column}
  \end{columns}
\end{frame}

\begin{frame}[fragile]{Condicionais}
  \begin{codigo}
    def avaliar(inflacao_12m, centro=3.0, tolerancia=1.5):
        if inflacao_12m > centro + tolerancia:
            return "acima do teto"
        elif inflacao_12m < centro - tolerancia:
            return "abaixo do piso"
        else:
            return "dentro do intervalo"

    avaliar(5.2)     # 'acima do teto'
  \end{codigo}
  \begin{itemize}
    \item A \alert{indentação} define os blocos (não há chaves nem \texttt{end})
    \item Comparações: \py{==  !=  <  <=  >  >=}; combinação: \py{and  or  not}
  \end{itemize}
\end{frame}

\begin{frame}[fragile]{Laços e funções}
  Inflação acumulada: $\displaystyle \pi^{\text{acum}} = \prod_{t=1}^{T}(1+\pi_t) - 1$
  \begin{codigo}
    def acumulada(taxas_pct):
        """Inflação acumulada (%) a partir de taxas mensais (%)."""
        fator = 1.0
        for pi in taxas_pct:          # percorre cada elemento
            fator = fator * (1 + pi / 100)
        return 100 * (fator - 1)

    acumulada([0.16, 0.43, 0.56, 0.43])   # 1.5889...
  \end{codigo}
  \begin{itemize}
    \item \py{def} cria uma função: \alert{entradas} $\to$ \alert{saída} (\py{return})
    \item A \emph{docstring} documenta o que a função faz
  \end{itemize}
\end{frame}

\begin{frame}{O laço passo a passo: simule o computador}
  \begin{columns}[c]
    \begin{column}{0.46\textwidth}
      \centering\small
      \begin{tabular}{@{}clc@{}}
        \toprule
        Volta & \py{pi} & \py{fator} depois da volta \\
        \midrule
        início & ---  & 1,0000 \\
        1      & 0,16 & 1,0016 \\
        2      & 0,43 & 1,0059 \\
        3      & 0,56 & 1,0115 \\
        4      & 0,43 & 1,0159 \\
        \bottomrule
      \end{tabular}
      \par\medskip
      \py{return}: $100 \times (1{,}0159 - 1) \approx 1{,}589$\%
    \end{column}
    \begin{column}{0.5\textwidth}
      \begin{intuicao}
        \small Cada volta do \py{for} pega o próximo elemento e
        \alert{atualiza} \py{fator}. Montar essa tabela com lápis é a melhor
        forma de entender (e de depurar) um laço.
      \end{intuicao}
      \begin{conexao}
        \small É a coluna de ``fator acumulado'' que você arrastaria numa
        planilha: cada linha da tabela é uma linha da planilha.
      \end{conexao}
    \end{column}
  \end{columns}
\end{frame}

\begin{frame}[fragile]{Teste rápido: o que o Python devolve?}
  \begin{codigo}
    ipca = [0.16, 0.43, 0.56, 0.43]
    ipca[-2:]        # (a)
    ipca * 2         # (b)
    avaliar(4.5)     # (c)  centro 3.0, tolerância 1.5
  \end{codigo}
  \begin{uncoverenv}<2>
  \begin{enumerate}[(a)]\small
    \item \py{[0.56, 0.43]}: os dois últimos elementos.
    \item \py{[0.16, 0.43, 0.56, 0.43, 0.16, 0.43, 0.56, 0.43]}: a lista é
          \alert{repetida}, não multiplicada! Para fazer conta com vetores
          precisamos do NumPy.
    \item \py{'dentro do intervalo'}: $4{,}5 > 4{,}5$ é falso. Atenção às
          fronteiras ($>$ \emph{versus} $\geq$).
  \end{enumerate}
  \end{uncoverenv}
\end{frame}

% =====================================================================
\section{NumPy: computação vetorizada}

\begin{frame}{Por que não basta usar listas?}
  \begin{columns}[c]
    \begin{column}{0.38\textwidth}
      Um \alert{array} NumPy é um bloco de memória com elementos do
      \textbf{mesmo tipo}. Operações sobre o array inteiro rodam em código
      compilado (C), não no interpretador Python.
      \medskip

      \textbf{Vetorizar} = trocar o laço explícito por uma operação sobre o array.
    \end{column}
    \begin{column}{0.6\textwidth}
      \includegraphics[width=\linewidth]{a02_laco_vs_vetor}
    \end{column}
  \end{columns}
\end{frame}

\begin{frame}[fragile]{Arrays: criação, forma e indexação}
  \begin{codigo}
    import numpy as np

    pi = np.array([0.16, 0.43, 0.56, 0.43])   # vetor, shape (4,)
    X = np.array([[1, 2.0], [1, 3.5], [1, 1.2]])  # matriz, shape (3, 2)

    X.shape          # (3, 2): 3 observações, 2 variáveis
    X[:, 1]          # segunda coluna -> array([2. , 3.5, 1.2])
    X[X[:, 1] > 1.5] # filtro booleano: linhas com x > 1.5

    np.zeros(5), np.ones((2, 3)), np.linspace(0, 1, 11)
  \end{codigo}
  \begin{itemize}
    \item Convenção do curso: \alert{linhas = observações}, \alert{colunas = variáveis}
  \end{itemize}
\end{frame}

\begin{frame}[fragile]{Eixos: resumir por coluna ou por linha?}
  \begin{columns}[c]
    \begin{column}{0.5\textwidth}
      \centering\small
      \begin{tabular}{@{}r|cc|c@{}}
        & \textbf{const.} & \textbf{renda} & \py{axis=1} \\
        \hline
        obs.~1 & 1 & 2,0 & 1,50 \\
        obs.~2 & 1 & 3,5 & 2,25 \\
        obs.~3 & 1 & 1,2 & 1,10 \\
        \hline
        \py{axis=0} & 1,00 & 2,23 & \\
      \end{tabular}
      \par\medskip
      \begin{codigo}
        X.mean(axis=0)  # [1.  , 2.23]
        X.mean(axis=1)  # [1.5 , 2.25, 1.1]
      \end{codigo}
    \end{column}
    \begin{column}{0.46\textwidth}
      \begin{intuicao}
        \small \py{axis=0} \alert{percorre as linhas} e devolve um número por
        coluna: a média de cada variável. É o que quase sempre queremos.
      \end{intuicao}
      \begin{atencao}
        \small \py{X.mean()} sem \py{axis} tira a média de \textbf{todos} os
        elementos, misturando variáveis.
      \end{atencao}
    \end{column}
  \end{columns}
\end{frame}

\begin{frame}[fragile]{Operações vetorizadas e \emph{broadcasting}}
  \begin{codigo}
    acum = 100 * (np.prod(1 + pi / 100) - 1)     # sem laço!

    salarios = np.array([2500., 4000., 7200.])  # nominais
    deflator = np.array([1.00, 1.05, 1.11])     # índice de preços
    reais = salarios / deflator                 # elemento a elemento

    X_pad = (X - X.mean(axis=0)) / X.std(axis=0) # (3,2) - (2,) -> (3,2)
  \end{codigo}
  \begin{intuicao}
    \small \emph{Broadcasting}: quando as formas diferem, o NumPy ``estica'' a
    menor ao longo dos eixos faltantes. Em \py{X - X.mean(axis=0)}, o vetor de
    médias (uma por coluna) é subtraído de cada linha.
  \end{intuicao}
\end{frame}

\begin{frame}[fragile]{Álgebra matricial: MQO em uma linha}
  \begin{columns}[T]
    \begin{column}{0.5\textwidth}
      \begin{conexao}
        \small Modelo $\vy = \mX\vbeta + \ve$. O estimador de MQO resolve
        \[ (\mX\T\mX)\,\hvbeta = \mX\T\vy. \]
      \end{conexao}
      \small Ao programar, \alert{resolva o sistema}; não calcule a inversa
      $(\mX\T\mX)^{-1}$ explicitamente (mais lento e numericamente instável).
    \end{column}
    \begin{column}{0.47\textwidth}
      \begin{codigo}
        rng = np.random.default_rng(42)
        n = 500
        x = rng.normal(size=n)
        X = np.column_stack(
              [np.ones(n), x])
        y = 1.0 + 2.0 * x \
            + rng.normal(size=n)

        beta = np.linalg.solve(
              X.T @ X, X.T @ y)
        # array([0.98, 2.03])
      \end{codigo}
    \end{column}
  \end{columns}
\end{frame}

\begin{frame}{Você já fazia isso: da planilha e do Stata ao NumPy}
  \centering\small
  \begin{tabular}{@{}>{\raggedright\arraybackslash}p{0.32\textwidth}
                    >{\raggedright\arraybackslash}p{0.2\textwidth}
                    >{\raggedright\arraybackslash}p{0.4\textwidth}@{}}
    \toprule
    Planilha & Stata & NumPy \\
    \midrule
    uma coluna de dados & uma variável & um array 1D (uma coluna de \py{X}) \\
    fórmula arrastada & \texttt{gen r = s / d} & \py{r = s / d} (vetorizado) \\
    \texttt{=MÉDIA(B2:B500)} & \texttt{summarize} & \py{X.mean(axis=0)} \\
    suplemento de regressão & \texttt{reg y x} & \py{np.linalg.solve(X.T @ X, X.T @ y)} \\
    planilha com cliques & \emph{do-file} & notebook ou script \\
    \bottomrule
  \end{tabular}
  \par\medskip\raggedright
  \begin{conexao}
    \small A lógica é a mesma; a diferença é que no NumPy \alert{escrevemos a
    conta}. Por isso, na Aula~4, trocaremos a fórmula do MQO por outro método
    de treino sem mudar de ferramenta.
  \end{conexao}
\end{frame}

\begin{frame}[fragile]{Teste rápido: formas e \emph{broadcasting}}
  \begin{codigo}
    X = np.ones((3, 2))
    X.mean(axis=0).shape             # (a)
    X + np.array([10, 20])           # (b)
    X + np.array([1, 2, 3])          # (c)
  \end{codigo}
  \begin{uncoverenv}<2>
  \begin{enumerate}[(a)]\small
    \item \py{(2,)}: uma média por coluna.
    \item Funciona, com forma \py{(3, 2)}: o vetor \py{[10, 20]} é somado a cada linha.
    \item \py{ValueError}: as formas \py{(3, 2)} e \py{(3,)} não combinam. O
          \emph{broadcasting} alinha as formas \alert{pela direita}: 2 contra 3.
  \end{enumerate}
  \end{uncoverenv}
\end{frame}

% =====================================================================
\section{JAX: NumPy que deriva e compila}

\begin{frame}{O que é o JAX?}
  Biblioteca do Google com a \alert{mesma interface do NumPy} e três
  transformações de funções:
  \begin{center}\small
  \begin{tabular}{@{}l>{\raggedright\arraybackslash}p{0.42\textwidth}l@{}}
    \toprule
    Transformação & O que faz & Onde usaremos \\
    \midrule
    \py{jax.grad} & derivada/gradiente \alert{exato} de uma função & treinar modelos (Aulas 4–5) \\
    \py{jax.jit}  & compila a função (XLA) para CPU, GPU ou TPU & acelerar o treino \\
    \py{jax.vmap} & vetoriza uma função escrita para um único ponto & avaliar em lotes \\
    \bottomrule
  \end{tabular}
  \end{center}
  \begin{atencao}
    \small Arrays JAX são \textbf{imutáveis}: em vez de \py{x[0] = 1}, escreva
    \py{x = x.at[0].set(1)}. Números aleatórios exigem uma \emph{chave} explícita.
  \end{atencao}
\end{frame}

\begin{frame}{Transformações se compõem}
  \centering
  \includegraphics[width=0.78\linewidth]{a02_jax_transformacoes}
  \par\medskip\raggedright
  Cada transformação \alert{recebe uma função e devolve outra função}. Por isso
  podemos encadeá-las: \py{jax.jit(jax.grad(perda))} é o motor do treinamento de
  modelos que construiremos nas Aulas~4 e~5.
\end{frame}

\begin{frame}[fragile]{Diferenciação automática: utilidade marginal}
  \begin{columns}[T]
    \begin{column}{0.46\textwidth}
      Utilidade CRRA: $u(c) = \dfrac{c^{1-\gamma}}{1-\gamma}$, \quad $u'(c) = c^{-\gamma}$
      \begin{codigo}
        import jax
        import jax.numpy as jnp

        gamma = 2.0
        def u(c):
            return c**(1-gamma) / (1-gamma)

        du = jax.grad(u)  # nova função!
        du(1.5)           # 0.444 = 1.5**-2
      \end{codigo}
    \end{column}
    \begin{column}{0.52\textwidth}
      \includegraphics[width=\linewidth]{a02_crra_grad}
      \small \py{jax.grad} não aproxima por diferenças finitas nem manipula
      símbolos: aplica a \alert{regra da cadeia} sobre as operações executadas.
    \end{column}
  \end{columns}
\end{frame}

\begin{frame}{Três formas de obter $u'(1{,}5)$}
  \begin{columns}[T]
    \begin{column}{0.48\textwidth}
      \textbf{1. À mão:} $u'(c) = c^{-\gamma} = 1{,}5^{-2} = 0{,}444444$.
      Exata, mas trabalhosa em modelos grandes.
      \smallskip

      \textbf{2. Diferenças finitas:} $\dfrac{u(c+h) - u(c)}{h}$
      \par\smallskip\centering\small
      \begin{tabular}{@{}lc@{}}
        \toprule
        $h$ & aproximação \\
        \midrule
        $10^{-2}$  & 0,441501 \\
        $10^{-6}$  & 0,444444 \\
        $10^{-13}$ & 0,444089 \\
        \bottomrule
      \end{tabular}
    \end{column}
    \begin{column}{0.48\textwidth}
      \textbf{3. Automática (\py{jax.grad}):} exata e sem derivar nada à mão.
      \smallskip
      \begin{atencao}
        \small $h$ grande erra pela aproximação; $h$ pequeno demais erra pelo
        arredondamento, o mesmo problema de \py{0.1 + 0.2}.
      \end{atencao}
      \begin{intuicao}
        \small Treinar exige o gradiente em milhares de parâmetros, a cada
        passo. Só a forma 3 escala.
      \end{intuicao}
    \end{column}
  \end{columns}
\end{frame}

\begin{frame}[fragile]{Prévia: o gradiente de uma função de perda}
  \begin{columns}[T]
    \begin{column}{0.5\textwidth}
      Erro quadrático médio como função dos parâmetros:
      \[ \loss(\vbeta) = \tfrac{1}{n}\norm{\vy - \mX\vbeta}^2 \]
      \[ \grad_{\vbeta}\loss = -\tfrac{2}{n}\mX\T(\vy - \mX\vbeta) \]
      \begin{conexao}
        \small Igualar esse gradiente a zero dá exatamente as equações normais do
        MQO. Na Aula~4 vamos \alert{caminhar} na direção oposta ao gradiente em vez
        de resolvê-las.
      \end{conexao}
    \end{column}
    \begin{column}{0.47\textwidth}
      \begin{codigo}
        def perda(beta, X, y):
            r = y - X @ beta
            return jnp.mean(r**2)

        g = jax.grad(perda)  # derivada em beta
        b0 = jnp.zeros(2)
        g(b0, X, y)

        # confere com a fórmula analítica:
        -2 / n * X.T @ (y - X @ b0)

        perda_rapida = jax.jit(perda)
      \end{codigo}
    \end{column}
  \end{columns}
\end{frame}

% =====================================================================
\section{Mãos à obra}

\demo[\nbdir/02_demo.ipynb]{Do laço ao gradiente:\\ IPCA, MQO e \texttt{jax.grad} ao vivo}

\lab[\nbdir/02_lab.ipynb]{Inflação acumulada em 12 meses\\ com dados do Banco Central}

\begin{frame}{Para fixar}
  \begin{itemize}
    \item Python: variáveis, tipos, coleções, \py{if}, \py{for}, \py{def}
    \item NumPy: arrays com \py{shape}; operações vetorizadas e
          \emph{broadcasting}; álgebra com \py{@} e \py{np.linalg.solve}
    \item JAX: interface do NumPy + \py{grad}, \py{jit}, \py{vmap}
    \item O gradiente de uma perda é o \alert{elo} entre econometria e treinamento de modelos
  \end{itemize}
  \medskip
  \textbf{Para aprofundar:} \textcite{downey2016} (cap.~1–10); documentação
  \emph{JAX — The Sharp Bits}; \textcite{baydin2018} sobre diferenciação automática;
  Lista~2 e notas de aula.
  \medskip

  \textbf{Próxima aula:} como saber se o modelo aprendeu ou apenas decorou?
\end{frame}

\begin{frame}{Referências}
  \nocite{harris2020,bradbury2018}
  \printbibliography[heading=none]
\end{frame}

\end{document}
